\section{Conformance}
\label{sec:conformance}

\subsection{KAT vector files}

\texttt{pqcrypt/conformance/vectors/}:

\begin{itemize}
\item \texttt{encrypt\_roundtrip.json} --- fixed-RNG-seed determinism. Each record gives a plaintext, a recipient public key, a sender RNG seed, and an expected envelope hex.
\item \texttt{decrypt.json} --- known envelopes $\to$ known plaintexts, exercising the decrypt path for suites \code{0x01} and \code{0x02}.
\item \texttt{peek\_header.json} --- header parsing without authentication.
\item \texttt{invalid.json} --- bad variant byte (covering 0x00, 0x01, 0x02, 0x03, 0x05--0xFF), bad suite byte, truncated \code{KEM\_CT}, bad AEAD tag, bad commit tag, bad AD.
\item \texttt{ad\_binding.json} --- \pqid{}-as-AD positive and negative cases.
\item \texttt{commitment.json} --- UtC commit tag KATs independent of the AEAD path.
\item A KyberSlash-class timing-regression vector (working name \texttt{kyberslash\_regression.json}) --- known-bad ML-KEM-1024 inputs that exercised KyberSlash 1/2 \citep{kyberslash}, asserting they reject without timing variance --- is \emph{planned} (backlog) and is \emph{not} yet part of the shipped vector set.
\item \texttt{min\_timestamp.json} --- \code{MIN\_TIMESTAMP\_MS} helper parity (mirrors \pqid{}'s and \pqhash{}'s analogous files).
\item \texttt{hybrid\_decrypt.json} --- KAT round-trip for hybrid suite \code{0x06} (4 records).
\item \texttt{hybrid\_ad\_binding.json} --- hybrid AD-binding positive and negative cases (4 records).
\item \texttt{hybrid\_peek\_header.json} --- hybrid header parse (1 record).
\end{itemize}

\subsection{Cross-language runner}

Vectors are consumed by the same harness used by \pqid{} and \pqhash{}: \texttt{tools/conformance/run.sh}. The script discovers per-language runners (\texttt{tools/conformance/python/}, \texttt{.../c/}, \texttt{.../go/}, \texttt{.../rust/}) and executes each against each vector file. Exit code 0 means every runner passed every vector. Language ports without their own \texttt{pqcrypt\_runner} fall back to the Python runner for PQCrypt vectors via the harness's \texttt{pqcrypt\_runner\_cmd} empty case (PQCK-00015).

\subsection{Hybrid suite \code{0x06} skip discipline}

\texttt{tools/conformance/run.sh} filters \texttt{hybrid\_*.json} vector files per language via \texttt{hybrid\_supported()}:

\begin{lstlisting}[language=bash]
hybrid_supported() {
    case "$1" in
        python) return 0 ;;   # PQCK-00021, v1.2
        rust)   return 0 ;;   # PQCK-00028, v1.3
        go)     return 0 ;;   # PQCK-00029, v1.3
        *)      return 1 ;;   # C/C++/TS/Java/CPython ext: v1.4 backlog
    esac
}
\end{lstlisting}

Languages without hybrid skip the \texttt{hybrid\_*.json} files entirely and MUST reject suite \code{0x06} as \texttt{UNKNOWN\_SUITE} for any non-hybrid envelope. When porting hybrid to a new language: (1) implement hybrid in the language, (2) update the language's conformance runner to accept 96-byte seeds (\cref{sec:hybrid}), (3) flip the language's \texttt{hybrid\_supported()} entry, (4) verify the language's PASS count rises by 9 (one per record across the three \texttt{hybrid\_*.json} files).

\subsection{Local invocation}

\begin{lstlisting}[language=bash]
bash tools/conformance/run.sh
PQTK_CONFORMANCE_SKIP="cpython java cpp typescript" bash tools/conformance/run.sh
\end{lstlisting}

Latest state (v1.3, 2026-05-21, Apple M4 Max):

\begin{center}
\small
\begin{tabular}{@{}lrrl@{}}
\toprule
\textbf{Language} & \textbf{PASS} & \textbf{FAIL} & \textbf{Notes} \\
\midrule
Python   & 265 & 0 & 256 standard + 9 hybrid \\
C        & 256 & 0 & hybrid skipped (v1.4 backlog) \\
Go       & 265 & 0 & 256 + 9 hybrid (PQCK-00029) \\
Rust     & 265 & 0 & 256 + 9 hybrid (PQCK-00028) \\
\bottomrule
\end{tabular}
\end{center}

C++, TypeScript, Java, and the CPython extension consume PQCrypt vectors via the Python runner fallback and skip \texttt{hybrid\_*.json}.

\subsection{Branch-protection gate}

The GitHub Actions workflow \texttt{.github/workflows/conformance.yml} runs the full harness across all three projects (\pqid, \pqhash, \pqcrypt). It is a required check for merging to \texttt{main}.

\subsection{Vector generation}

All vectors are generated against the Python reference (\texttt{pqcrypt/python/pqcrypt.py}). The generator script lives at \texttt{pqcrypt/conformance/scripts/generate.py}. Vectors are committed verbatim: re-running the generator on the same code MUST produce identical JSON. If a code change produces different vectors, that is a spec change and requires a release-notes entry plus updates to every language implementation.
